\SourceLeaf{284}%
\[
A\int\frac{\cos y}{m\cos\psi}\mathrm{d}x
\]
von \(x=0\) bis \(x=h\). Man darf darin aber stets \(\cos y\) gleich eins setzen, \(\cos\psi=1\) wählen und für \(m\) den Wert verwenden, der auf der Linie \(M\) oder für \(y=0\) gilt. Damit wird das Integral
\[
\begin{aligned}
&=A\int\frac{\mathrm{d}x}{m^0(1+\mu x+\mu'xx+\text{usw.})}\\
&=\frac{A}{m^0}(h-\tfrac{1}{2}\mu hh+(\tfrac{1}{3}\mu\mu-\tfrac{1}{3}\mu')h^3-\text{usw.})
\end{aligned}
\]
Es genügt stets, dafür den bis zur Ordnung \(hh\) damit übereinstimmenden Wert
\[
\frac{Ah}{\BareRoot{m^0m'}}
\]
zu verwenden.

\Article{13}
Die Größen \(l^0,l'\) werden auf folgende Weise bestimmt. Es sei \(\chi\) der Winkel, den der Großkreisbogen \(N\) an einer beliebigen Stelle in Richtung wachsender \(x\) mit dem von Norden nach Süden gerichteten Meridian bildet. Der Winkel wird vom Meridian zum Bogen von links nach rechts gezählt. Ferner sei \(S\) die Breite an jener Stelle und \(T\) die von einem beliebigen Meridian aus ostwärts gezählte Länge. An dieser Stelle gilt dann
\[
\begin{aligned}
\mathrm{d}S&=-\cos\chi\mathord{.}\mathrm{d}x+\sin\chi\mathord{.}\mathrm{d}y\\
\mathrm{d}T&=-\frac{\sin\chi}{\cos S}\mathrm{d}x-\frac{\cos\chi}{\cos S}\mathrm{d}y
\end{aligned}
\]
und damit ergibt sich der partielle Differentialquotient
\[
\frac{\mathrm{d}m}{m\mathrm{d}y}=\sin\chi\mathord{.}\frac{\mathrm{d}m}{m\mathrm{d}S}-\frac{\cos\chi}{\cos S}\mathord{.}\frac{\mathrm{d}m}{m\mathrm{d}T}
\]

Bei unserer konformen Abbildung ist \(m\) von der Länge unabhängig, das heißt, es gilt \(\tfrac{\mathrm{d}m}{m\mathrm{d}T}=0\). Damit erhält man
\[
l=\sin\chi\mathord{.}\frac{\mathrm{d}m}{m\mathrm{d}S}
\]

Bezeichnet man die Werte von \(\chi\) in den Punkten \(F\) und \(G\) mit \(V^0\) und \(180^0+V'\) (sodass nach üblichem Sprachgebrauch \(V^0\) das Azimut des Großkreisbogens \(FG\) in \(F\) und \(V'\) das Azimut des Großkreisbogens
